% (c) Nikita Lisitsa, lisyarus@gmail.com, 2026

\documentclass[10pt]{beamer}

\usepackage[T2A]{fontenc}
\usepackage[russian]{babel}
\usepackage{minted}

\usepackage[most]{tcolorbox}
\tcbuselibrary{minted}

\usepackage{graphicx}
\graphicspath{ {./images/} }

\usepackage{adjustbox}

\usepackage{color}
\usepackage{soul}
\usepackage{multicol}
\usepackage{multirow}
\usepackage{hyperref}
\usepackage{tikz}
\usetikzlibrary{arrows.meta,positioning,calc}

\usetheme{metropolis}

\definecolor{red}{rgb}{1,0,0}
\definecolor{codebg}{RGB}{29,35,49}
\setminted{fontsize=\footnotesize}

\makeatletter
\newcommand{\slideimage}[1]{
  \begin{figure}
    \begin{adjustbox}{width=\textwidth, totalheight=\textheight-2\baselineskip-2\baselineskip,keepaspectratio}
      \includegraphics{#1}
    \end{adjustbox}
  \end{figure}
}
\makeatother

\title{Язык программирования C}
\subtitle{Лекция 2: встроенные типы, массивы, указатели}
\author{Лисица Никита}
\date{2026}

\setbeamertemplate{footline}[frame number]

\usemintedstyle{tango}

\begin{document}

\frame{\titlepage}

\begin{frame}[fragile]
\frametitle{Встроенные (базовые) типы данных}
На большинстве современных платформ:
\begin{center}
  \begin{tabular}{ c c c }
    \hline
    Тип & Размер в байтах & Диапазон значений \\
    \hline
    \mintinline{cpp}|char| & 1 & \(-2^7\dots 2^7-1\) \\
    \mintinline{cpp}|short| & 2 & \(-2^{15}\dots 2^{15}-1\) \\
    \mintinline{cpp}|int| & 4 & \(-2^{31}\dots 2^{31}-1\) \\
    \mintinline{cpp}|long| & 4/8 & \(-2^{31/63}\dots 2^{31/63}-1\) \\
    \mintinline{cpp}|long long| & 8 & \(-2^{63}\dots 2^{63}-1\) \\
    \mintinline{cpp}|unsigned char| & 1 & \(0\dots 2^8-1\) \\
    \mintinline{cpp}|unsigned short| & 2 & \(0\dots 2^{16}-1\) \\
    \mintinline{cpp}|unsigned int| & 4 & \(0\dots 2^{32}-1\) \\
    \mintinline{cpp}|unsigned long| & 4/8 & \(0\dots 2^{32/64}-1\) \\
    \mintinline{cpp}|unsigned long long| & 8 & \(0\dots 2^{64}-1\) \\
    \hline
    \mintinline{cpp}|float| & 4 & \(\pm1.4\times 10^{-45}\dots 3.4\times 10^{38}\) \\
    \mintinline{cpp}|double| & 8 & \(\pm4.9\times 10^{-324}\dots 1.8\times 10^{308}\) \\
    \hline
  \end{tabular}
\end{center}
\pause
\begin{itemize}
\item Почему `на большинстве'?
\end{itemize}
\end{frame}

\begin{frame}[fragile]
\frametitle{Встроенные (базовые) типы данных}
\begin{itemize}
\item Стандарт языка C гарантирует только \textbf{минимальный размер} встроенных типов и неравенства на их порядок
\pause
\item \mintinline{cpp}|char| -- всегда 1 байт (почти всегда 8 бит)
\pause
\item \mintinline{cpp}|int| -- естественный для CPU тип (\(\approx\) размер регистров), как минимум 16 бит (обычно 32)
\pause
\item \mintinline{cpp}|sizeof(short)| \(\leq\) \mintinline{cpp}|sizeof(int)| \(\leq\) \mintinline{cpp}|sizeof(long)|
\pause
\item \mintinline{cpp}|sizeof| -- встроенный в язык оператор (не функция!), в момент компиляции его вызов заменяется на размер типа или значения
\end{itemize}
\end{frame}

\begin{frame}[fragile]
\frametitle{Изучить (вспомнить) самостоятельно}
\begin{itemize}
\item Представление отрицательных чисел в дополнительном коде
\pause
\item Представление чисел с плавающей точкой
\pause
\item Приведение типов (явное, неявное)
\pause
\item Что-то из этого будет в письменном тесте :)
\end{itemize}
\end{frame}

\begin{frame}[fragile]
\frametitle{Массивы}
\begin{itemize}
\item Последовательность объектов одного типа, лежащих в непрерывном участке памяти
\item Размер должен быть известен \textbf{на момент компиляции}!
\end{itemize}
\usemintedstyle{lightbulb}
\begin{tcblisting}{listing engine=minted, minted language=cpp, minted options={fontsize=\scriptsize}, listing only, colback=codebg, colframe=codebg, rounded corners}
// Массив из 10-ти элементов типа int
// Размер в памяти: 10 * sizeof(int)
int array[10];

// Инициализация:
int array[10] = {0, 1, 2, 3, 4, 5, 6, 7, 8, 9};

// Заполнить нулями:
int array[10] = {0};

// Вывести размер из инициализации:
int array[] = {0, 1, 2, 3, 4, 5, 6, 7, 8, 9};
\end{tcblisting}
\end{frame}

\begin{frame}[fragile]
\frametitle{Массивы}
\usemintedstyle{lightbulb}
\begin{tcblisting}{listing engine=minted, minted language=cpp, minted options={fontsize=\scriptsize}, listing only, colback=codebg, colframe=codebg, rounded corners}
char array[10];

// Прочитать/записать значение по индексу:
array[0] = 10;
array[1] = array[0] + 5;

// Узнать размер:
int size = sizeof(array) / sizeof(array[0]);
\end{tcblisting}
\end{frame}

\begin{frame}[fragile]
\frametitle{Массивы char}
\usemintedstyle{lightbulb}
\begin{tcblisting}{listing engine=minted, minted language=cpp, minted options={fontsize=\scriptsize}, listing only, colback=codebg, colframe=codebg, rounded corners}
// Массив из 10-ти элементов типа char
// Размер в памяти: 10 * sizeof(int)
char array[10];

char array[10] = {0, 1, 2, 3, 4, 5, 6, 7, 8, 9};

char array[] = {'H', 'e', 'l', 'l', 'o', ',',
  ' ', 'w', 'o', 'r', 'l', 'd', '!', '\0'};

// То же самое через строковой литерал
// 14 элементов - включает терминирующий '\0'
char array[] = "Hello, world!";

// NB: вот такой синтаксис - абсолютно другая операция!
char *array = "Hello, world!";
\end{tcblisting}
\end{frame}

\begin{frame}[fragile]
\frametitle{Многомерные массивы}
\usemintedstyle{lightbulb}
\begin{tcblisting}{listing engine=minted, minted language=cpp, minted options={fontsize=\scriptsize}, listing only, colback=codebg, colframe=codebg, rounded corners}
// Массив из 3-х элементов, каждый - массив из 4-х int-ов
// Размер в памяти: 3 * 4 * sizeof(int)
int array[3][4];

int array[2][2] = {
  {1, 2},
  {3, 4},
};
\end{tcblisting}
\end{frame}

\begin{frame}[fragile]
\frametitle{Выход за пределы массива}
\begin{itemize}
\item Для массива \mintinline{cpp}|int array[N];| обращение по индексу \(i<0\) или \(i\geq N\) -- \textbf{undefined behavior}
\pause
\item Что произойдёт на практике, если написать \mintinline{cpp}|array[-1]|?
\pause
\item На современных компиляторах будет предупреждение/ошибка
\pause
\item Если индекс неизвестен в момент компиляции (переменная/аргумент функции/т.п.), процессор попытается прочитать память за пределами массива
\pause
\item Если это память того же процесса -- ошибки не будет, но будет прочитано `мусорное' значение
\pause
\item Если это память другого процесса -- access violation / segmentation fault
\end{itemize}
\end{frame}

\begin{frame}[fragile]
\frametitle{Указатели}
\begin{itemize}
\item Указатель (pointer) -- объект, указывающий/ссылающийся на другой объект
\pause
\item На практике на большинстве платформ указатель -- адрес в памяти (число), номер первого байта объекта
\pause
\item \textbf{N.B.:} это адрес не в \textit{физической памяти}, а в \textit{виртуальной}, но это не принципиально для понимания
\pause
\item \mintinline{cpp}|int *p;| -- указатель на кусок памяти, в котором хранится число типа \mintinline{cpp}|int|
\end{itemize}
\end{frame}

\begin{frame}[fragile]
\frametitle{Указатели}
\usemintedstyle{lightbulb}
\begin{tcblisting}{listing engine=minted, minted language=cpp, minted options={fontsize=\scriptsize}, listing only, colback=codebg, colframe=codebg, rounded corners}
int a = 42;

// Взятие адреса переменной:
int *p = &a;

// "Разыменование" указателя
// Прочитать значение по адресу
int b = *p;

// "Разыменование" указателя
// Записать значение по адресу
*p = 67;

// Напечатать указатель
printf("%p", p);
\end{tcblisting}
\end{frame}

\begin{frame}[fragile]
\frametitle{Размер указателя}
\begin{itemize}
\item Стандарт C не регламентирует размер самого указателя, но он одинаков для всех типов указателей (кроме указателей на функции)
\pause
\item На практике размер указателя зависит от битности системы: 4 байта на 32-битной (x86), 8 байт на 64-битной (x86\_64, ARM64)
\end{itemize}
\end{frame}

\begin{frame}[fragile]
\frametitle{Размер указателя}
\begin{itemize}
\item Не путайте размер указателя и размер объекта, не который ссылается указатель!
\end{itemize}
\usemintedstyle{lightbulb}
\begin{tcblisting}{listing engine=minted, minted language=cpp, minted options={fontsize=\scriptsize}, listing only, colback=codebg, colframe=codebg, rounded corners}
char c = 'A';
int i = 10;
double d = 3.1415;

// На x86\_64
sizeof(c) == 1;
sizeof(i) == 4;
sizeof(d) == 8;
sizeof(&c) == 8;
sizeof(&i) == 8;
sizeof(&d) == 8;
\end{tcblisting}
\end{frame}

\begin{frame}[fragile]
\frametitle{Арифметика указателей}
\begin{itemize}
\item К указателям можно прибавлять/вычитать число
\pause
\item Семантика такой операции -- перейти к следующему элементу того же типа в непрерывном участке памяти/массиве
\pause
\item \(\Longrightarrow\) К адресу, хранящемуся в указателе, прибавляется число, умноженное на \mintinline{cpp}|sizeof(*p)|
\end{itemize}
\end{frame}

\begin{frame}[fragile]
\frametitle{Арифметика указателей}
\usemintedstyle{lightbulb}
\begin{tcblisting}{listing engine=minted, minted language=cpp, minted options={fontsize=\scriptsize}, listing only, colback=codebg, colframe=codebg, rounded corners}
char array_c[] = "MKN rules";
int array_i[10];
double array_d[10];

char *pc = &array_c[0];
int *pi = &array_i[0];
double *pd = &array_d[0];

// К фактическому адресу прибавится sizeof(char) == 1 байт
pc += 1;

// К фактическому адресу прибавится sizeof(int) == 4 байт
pi += 1;

// К фактическому адресу прибавится sizeof(double) == 8 байт
pd += 1;
\end{tcblisting}
\end{frame}

\begin{frame}[fragile]
\frametitle{Арифметика указателей}
\begin{itemize}
\item Аналогично, вычитание указателей даёт количество элементов между ними (а не байт!)
\usemintedstyle{lightbulb}
\begin{tcblisting}{listing engine=minted, minted language=cpp, minted options={fontsize=\scriptsize}, listing only, colback=codebg, colframe=codebg, rounded corners}
int array[10];

int *p1 = &array[3];
int *p2 = &array[7];

// delta == 4
int delta = p2 - p1;
\end{tcblisting}
\end{itemize}
\end{frame}

\begin{frame}[fragile]
\frametitle{Указатели и массивы}
\begin{itemize}
\item Массив автоматически приводится к указателю на первый элемент массива:
\usemintedstyle{lightbulb}
\begin{tcblisting}{listing engine=minted, minted language=cpp, minted options={fontsize=\scriptsize}, listing only, colback=codebg, colframe=codebg, rounded corners}
int array[10];
// То же самое, что p = &array[0]
int *p = array;

// Ошибка компиляции: массиву нельзя присвоить новый адрес
array = p;
\end{tcblisting}
\pause
\item Через указатель тоже можно обращаться по индексу;
\usemintedstyle{lightbulb}
\begin{tcblisting}{listing engine=minted, minted language=cpp, minted options={fontsize=\scriptsize}, listing only, colback=codebg, colframe=codebg, rounded corners}
int *p = ...;

p[2] = p[1] + p[0];
\end{tcblisting}
\pause
\usemintedstyle{tango}
\item \mintinline{cpp}|p[i]| эквивалентно \mintinline{cpp}|*(p + i)|
\end{itemize}
\end{frame}

\begin{frame}[fragile]
\frametitle{Указатели и массивы}
\begin{itemize}
\item \mintinline{cpp}|p[i]| эквивалентно \mintinline{cpp}|*(p + i)|
\pause
\item Сложение коммутативно: \mintinline{cpp}|*(p + i) == *(i + p)|
\pause
\item \(\Longrightarrow\) Такой синтаксис тоже работает: \mintinline{cpp}|i[p]|
\pause
\item Не стоит так делать в реальном коде :)
\end{itemize}
\end{frame}

\begin{frame}[fragile]
\frametitle{Приведение типов указателей}
\begin{itemize}
\item В языке C любые типы указателей приводятся друг к другу:
\usemintedstyle{lightbulb}
\begin{tcblisting}{listing engine=minted, minted language=cpp, minted options={fontsize=\scriptsize}, listing only, colback=codebg, colframe=codebg, rounded corners}
int  *pi = ...;
char *pc = ...;

// ОК в C, ошибка в C++
pi = pc;
\end{tcblisting}
\pause
\item В C++ потребуется явное приведение типа:
\usemintedstyle{lightbulb}
\begin{tcblisting}{listing engine=minted, minted language=cpp, minted options={fontsize=\scriptsize}, listing only, colback=codebg, colframe=codebg, rounded corners}
pi = (int*) pc;
\end{tcblisting}
\end{itemize}
\end{frame}

\begin{frame}[fragile]
\frametitle{Применение указателей}
\begin{itemize}
\item Хотим написать функцию, меняющую местами два значения
\usemintedstyle{lightbulb}
\begin{tcblisting}{listing engine=minted, minted language=cpp, minted options={fontsize=\scriptsize}, listing only, colback=codebg, colframe=codebg, rounded corners}
void swap(int a, int b) {
  int temp = a;
  a = b;
  b = temp;
}

int main() {
  int x = 3, y = 4;
  swap(x, y);
  return 0;
}
\end{tcblisting}
\pause
\usemintedstyle{tango}
\item Функция \mintinline{cpp}|swap| поменяет местами \textbf{копии} аргументов, а не сами переменные \mintinline{cpp}|x,y|
\end{itemize}
\end{frame}

\begin{frame}[fragile]
\frametitle{Применение указателей}
\begin{itemize}
\item Решение:
\usemintedstyle{lightbulb}
\begin{tcblisting}{listing engine=minted, minted language=cpp, minted options={fontsize=\scriptsize}, listing only, colback=codebg, colframe=codebg, rounded corners}
void swap(int *a, int *b) {
  int temp = *a;
  *a = *b;
  *b = temp;
}

int main() {
  int x = 3, y = 4;
  swap(&x, &y);
  return 0;
}
\end{tcblisting}
\pause
\usemintedstyle{tango}
\item В функцию \mintinline{cpp}|swap| передаются \textbf{копии адресов} переменных \mintinline{cpp}|x,y| \(\Longrightarrow\) функция может поменять значения самих переменных
\end{itemize}
\end{frame}

\begin{frame}[fragile]
\frametitle{Применение указателей}
\begin{itemize}
\item Хотим передать в функцию большой объект и избежать лишнего копирования:
\usemintedstyle{lightbulb}
\begin{tcblisting}{listing engine=minted, minted language=cpp, minted options={fontsize=\scriptsize}, listing only, colback=codebg, colframe=codebg, rounded corners}
big_type big_object = ...;
foo(&big_object);
\end{tcblisting}
\pause
\begin{tcblisting}{listing engine=minted, minted language=cpp, minted options={fontsize=\scriptsize}, listing only, colback=codebg, colframe=codebg, rounded corners}
char str[] = "Lorem ipsum dolor sit amet, ...";
int size = strlen(str);
\end{tcblisting}
\end{itemize}
\end{frame}

\begin{frame}[fragile]
\frametitle{Применение указателей}
\begin{itemize}
\item Хотим пройтись по большому массиву:
\end{itemize}
\begin{multicols}{2}
\usemintedstyle{lightbulb}
\begin{tcblisting}{listing engine=minted, minted language=cpp, minted options={fontsize=\scriptsize}, listing only, colback=codebg, colframe=codebg, rounded corners}
int strlen (char *ptr) {
  int len = 0;
  while (ptr[len] != '\0') {
    ++len;
  }
  return len;
}
\end{tcblisting}
\columnbreak
\begin{tcblisting}{listing engine=minted, minted language=cpp, minted options={fontsize=\scriptsize}, listing only, colback=codebg, colframe=codebg, rounded corners}
int strlen (char *ptr) {
  char *p = ptr;
  while (*p != '\0') {
    ++p;
  }
  return p - ptr;
}
\end{tcblisting}
\end{multicols}
\pause
\begin{itemize}
\usemintedstyle{tango}
\item \mintinline{cpp}|ptr[len]| раскрывается в \mintinline{cpp}|*(ptr + len)| \(\Longrightarrow\) лишнее сложение в сравнении с правым вариантом
\end{itemize}
\end{frame}

\end{document}
